Der Auftrieb ist die Gewichtskraft des verdrängten Wassers, $F_\text{A} = \rho_\text{W}\,V\,g$, und hängt damit vom Volumen $V$ des Körpers ab. Beim Ausatmen wird die Lunge kleiner, das Volumen des Körpers nimmt um einige Liter ab, während seine Masse und damit die Gewichtskraft praktisch gleich bleiben. Der Auftrieb wird \result{\text{kleiner als die Gewichtskraft}}, man sinkt. Anders gesagt: Die mittlere Dichte des Körpers steigt beim Ausatmen über die Dichte des Wassers.
